\documentclass[10pt,a4paper]{amsart}
\usepackage[margin=19mm]{geometry}
\usepackage{amsmath,amssymb,mathtools}
\usepackage[scr=rsfs,cal=euler]{mathalfa}
\usepackage{xfrac}
\usepackage[unicode,hidelinks,bookmarks=false]{hyperref}
\newcommand{\ie}{i.e.\ }
\newcommand{\cf}{cf.\ }
\newcommand{\resp}{respectively\ }
\newcommand{\Subsection}{\subsection}
\providecommand{\nobreakdash}{\nobreak\hbox{-}}
\setlength{\emergencystretch}{2em}
\multlinegap0pt
\usepackage{bm}
\usepackage{amssymb}
\usepackage{tikz}
\usepackage{enumerate}
\usepackage[shortlabels]{enumitem}
\usepackage{csquotes}
\usepackage{float}
\usepackage{xcolor}
\newtheorem{lemma}{Lemma}
\newtheorem{claim}{Claim}
\newtheorem{poc}{Poc}
\newcommand{\R}{\mathbb{R}}
\newcommand*{\abs}[1]{\lvert#1\rvert}
\title{Tverberg's theorem for unions of convex sets: Sharp bounds and colored extensions}
\author{Gennian Ge, Yang Shu, Zixiang Xu}
\date{Source: arXiv:2607.12449v2 (CC BY 4.0)}
\begin{document}
\maketitle

\noindent\emph{Source and licence.} Tverberg's theorem for unions of convex sets: Sharp bounds and colored extensions. Version arXiv:2607.12449v2 (CC BY 4.0), distributed by arXiv under the Creative Commons Attribution 4.0 International licence, http://creativecommons.org/licenses/by/4.0/, as recorded in the arXiv licence metadata for this exact version and shown on the abstract page https://arxiv.org/abs/2607.12449v2. Full licence notice: \texttt{licenses/arxiv-2607.12449-CC-BY-4.0.txt}.\par\smallskip

\noindent\textit{Editorial scope.} Counted unit: the article's lemma lem:two-part-amplification (source0.tex lines 579--581). The article's complete original proof of it is delivered with it (source0.tex lines 583--635, one \texttt{proof} environment of 53 lines). No mathematics is written, repaired or re-derived: the statement and the proof are the article's own lines, in the article's own notation, and no retained line is substituted or re-flowed.  No further source-local context is needed: every cross-reference of every retained line resolves inside the two delivered spans.  Nothing was omitted from any retained span, so the package carries no omission record.  No retained line contains a citation, so the wrapper carries no bibliography and no bibliography entry.  No retained line contains a figure environment, an image-inclusion command or drawing code, so no image is delivered and the package declares no asset dependency.  Editorial formatting only, in addition: an AMS \texttt{amsart} preamble that restates the article's own declarations and macros used by the retained lines, namely the environments \texttt{lemma}, \texttt{claim}, \texttt{poc} on separate counters, together with the standard packages those lines need.  The retained environments print in this document as Lemma 1, Claim 1, Poc 1 in order of appearance, whereas the article prints the same items with the numbering of its own full document.  The complete native source of the article as distributed in the arXiv e-print for this exact version is bundled unchanged as \texttt{sources/205114/source0.tex}.
\begin{lemma}\label{lem:two-part-amplification}
Let \(\ell\ge2\), and let \(X\subseteq\R^{\ell}\) be finite. Suppose that every subset of \(X\) is the intersection of \(X\) with a union of at most \(s\) closed halfspaces, each having a normal vector whose coordinates are strictly positive. For every integer \(t\ge1\), there is a set \(X'\subseteq\R^{\ell}\) of \(t\abs{X}\) points having property \(\mathsf{B}(s,t)\).
\end{lemma}

\begin{proof}[Proof of Lemma~\ref{lem:two-part-amplification}]
The geometric idea is to anchor \(t\) tiny affine copies of \(X\) at distinct vertices of a polytope. For a halfspace that realizes a prescribed trace on one copy, we transform its normal vector into the interior of the normal cone at that vertex. The transformed halfspace cuts the designated copy exactly as prescribed and contains every other copy. We can therefore group the first color according to the original \(s\) halfspaces and group the second color according to the \(t\) host vertices.

Write \(h(\boldsymbol{u},c)=\{\boldsymbol{x}:\langle\boldsymbol{u},\boldsymbol{x}\rangle\le c\}\). For every subset \(Y\subseteq X\), choose one representation of \(Y\) by at most \(s\) halfspaces as in the assumption. If fewer than \(s\) are used, add copies of a fixed halfspace with coordinatewise positive normal vector that contains no point of \(X\). We may therefore assume that exactly \(s\) halfspaces have been chosen for every \(Y\). Since \(X\) is finite, only finitely many pairs \((\boldsymbol{u},c)\) occur in all these chosen representations.

To implement this idea, choose a full-dimensional convex polytope \(Q\subseteq\R^{\ell}\) with distinct vertices \(\boldsymbol{v}_{1},\ldots,\boldsymbol{v}_{t}\). For each \(j\in[t]\), define
\[
N_{Q}(\boldsymbol{v}_{j})=\{\boldsymbol{w}\in\R^{\ell}:\langle\boldsymbol{w},\boldsymbol{z}-\boldsymbol{v}_{j}\rangle\le0\text{ for every }\boldsymbol{z}\in Q\}.
\]
This is the set of vectors whose corresponding linear functions are maximized over \(Q\) at \(\boldsymbol{v}_{j}\). Since \(\boldsymbol{v}_{j}\) is a vertex, the outward normal vectors of the facets containing it span \(\R^{\ell}\), their positive combinations show that \(N_{Q}(\boldsymbol{v}_{j})\) has nonempty interior. Choose linearly independent vectors \(\boldsymbol{w}_{j,1},\ldots,\boldsymbol{w}_{j,\ell}\) in this interior, and let \(L_{j}\) be the invertible linear map satisfying \(L_{j}\boldsymbol{e}_{b}=\boldsymbol{w}_{j,b}\) for every \(b\in[\ell]\). If every coordinate of \(\boldsymbol{u}\) is positive, then
\[
L_{j}\boldsymbol{u}=\sum_{b=1}^{\ell}u_{b}\boldsymbol{w}_{j,b}\in\operatorname{int}N_{Q}(\boldsymbol{v}_{j}).
\]
Here, \(\operatorname{int}N_{Q}(\boldsymbol{v}_{j})\) denotes the
interior of the normal cone in \(\R^{\ell}\). In particular, \(\boldsymbol{v}_{j}\) is the unique vertex of \(Q\) at which the linear function with normal vector \(L_{j}\boldsymbol{u}\) is maximized.

For \(\varepsilon>0\), put
\[
X^{(j)}=\{\boldsymbol{v}_{j}+\varepsilon L_{j}^{-\top}\boldsymbol{x}:\boldsymbol{x}\in X\}.
\]
Here \(L_{j}^{-\top}\) denotes the inverse transpose of \(L_{j}\).
For every chosen halfspace \(h(\boldsymbol{u},c)\), define
\[
\widehat H_{j}(\boldsymbol{u},c)=\{\boldsymbol{y}\in\R^{\ell}:\langle L_{j}\boldsymbol{u},\boldsymbol{y}-\boldsymbol{v}_{j}\rangle\le\varepsilon c\}.
\]
If \(\boldsymbol{y}=\boldsymbol{v}_{j}+\varepsilon L_{j}^{-\top}\boldsymbol{x}\), then
\[
\langle L_{j}\boldsymbol{u},\boldsymbol{y}-\boldsymbol{v}_{j}\rangle=\varepsilon\langle\boldsymbol{u},\boldsymbol{x}\rangle.
\]
Thus \(\widehat H_{j}(\boldsymbol{u},c)\) has exactly the same intersection with \(X^{(j)}\) as \(h(\boldsymbol{u},c)\) has with \(X\), under the affine identification between these two sets.

The same halfspace also contains every other copy when \(\varepsilon\) is sufficiently small. Indeed, if \(b\ne j\), then \(L_{j}\boldsymbol{u}\in\operatorname{int}N_{Q}(\boldsymbol{v}_{j})\) gives
\(
\langle L_{j}\boldsymbol{u},\boldsymbol{v}_{b}-\boldsymbol{v}_{j}\rangle<0.
\)
Consequently, for every \(\boldsymbol{x}\in X\), one has \(\boldsymbol{v}_{b}+\varepsilon L_{b}^{-\top}\boldsymbol{x}\in\widehat H_{j}(\boldsymbol{u},c)\) once \(\varepsilon>0\) is small enough. There are only finitely many vertices, points, and chosen halfspaces, so one value of \(\varepsilon\) works simultaneously for all of them. We also take it small enough that the sets \(X^{(1)},\ldots,X^{(t)}\) are pairwise disjoint, and put \(X'=X^{(1)}\sqcup\cdots\sqcup X^{(t)}\).

Fix a map \(\chi:X'\to[2]\), and put \(A=\chi^{-1}(1)\) and \(B=\chi^{-1}(2)\). For each \(j\in[t]\), pull the set \(A\cap X^{(j)}\) back to a subset of \(X\), and use its chosen representation by \(s\) halfspaces. Let \(\widehat H_{j,1},\ldots,\widehat H_{j,s}\) be the corresponding halfspaces in \(\R^{\ell}\). Their union contains exactly the points of \(A\cap X^{(j)}\) inside the \(j\)-th copy. Assign each such point to one halfspace that contains it, and let \(A_{j,a}\) be the set assigned to \(\widehat H_{j,a}\). Define \(A_{a}=\bigcup_{j=1}^{t}A_{j,a}\) for \(a\in[s]\), and define \(B_{j}=B\cap X^{(j)}\) for \(j\in[t]\). These sets partition \(A\) and \(B\), respectively.

The following claim verifies the required component-by-component disjointness.

\begin{claim}\label{claim:two-part-disjoint}
For every \(a\in[s]\) and \(j\in[t]\), one has \(\operatorname{conv}(A_{a})\cap\operatorname{conv}(B_{j})=\emptyset\).
\end{claim}

\begin{poc}
By construction, \(A_{j,a}\subseteq\widehat H_{j,a}\). Moreover, every copy \(X^{(b)}\) with \(b\ne j\) is contained in \(\widehat H_{j,a}\). Hence \(A_{a}\subseteq\widehat H_{j,a}\), and the convexity of this halfspace gives \(\operatorname{conv}(A_{a})\subseteq\widehat H_{j,a}\).

Inside \(X^{(j)}\), the union of \(\widehat H_{j,1},\ldots,\widehat H_{j,s}\) contains exactly the points of \(A\cap X^{(j)}\). Therefore every point of \(B_{j}\) lies in the open halfspace \(\R^{\ell}\setminus\widehat H_{j,a}\). This open halfspace is convex, so it also contains \(\operatorname{conv}(B_{j})\). The two convex hulls are therefore disjoint.
\end{poc}

Claim \ref{claim:two-part-disjoint} says exactly that the partitions \(A=A_{1}\sqcup\cdots\sqcup A_{s}\) and \(B=B_{1}\sqcup\cdots\sqcup B_{t}\) satisfy property \(\mathsf{B}(s,t)\).
\end{proof}

\end{document}
